% from ega0-20-fr.tex
\oldpage[0\textsubscript{IV-1}]{133}
\begin{env}[20.5.11]
\label{0.20.5.11-fr}
Soient $A$ un anneau, $B$ une $A$-algèbre, $\mathfrak{R}$ un idéal de $B$,
$C$ la $A$-algèbre quotient $B/\mathfrak{R}$, et considérons
l'homomorphisme composé de $A$-modules
\begin{equation}
\tag{20.5.11.1}
\label{0.20.5.11.1-fr}
  \delta:\mathfrak{R}\longrightarrow B\xrightarrow{d_{B/A}}\Omega^1_{B/A}
\end{equation}
où la première flèche est l'injection canonique; comme
$d(xy)=xdy+ydx$, on voit que l'on a
$\delta(\mathfrak{R}^2)\subset\mathfrak{R}.\Omega^1_{B/A}$, d'où, par
passage aux quotients, un homomorphisme de $A$-modules
\begin{equation}
\tag{20.5.11.2}
\label{0.20.5.11.2-fr}
  \delta_{C/B/A}:\mathfrak{R}/\mathfrak{R}^2\longrightarrow
  \Omega^1_{B/A}\otimes_B C=
  \Omega^1_{B/A}/(\mathfrak{R}.\Omega^1_{B/A}).
\end{equation}
Mais en fait, $\delta_{C/B/A}$ est un homomorphisme de $C$-modules, car
pour $x\in B$ et $y\in\mathfrak{R}$, on a
$ydx\in\mathfrak{R}.\Omega^1_{B/A}$, donc
$d(xy)\equiv xdy\pmod{\mathfrak{R}.\Omega^1_{B/A}}$, ce qui prouve
d'abord que (20.5.11.2) est un homomorphisme de $B$-modules, et comme
$\mathfrak{R}$ annule les deux membres, cela établit notre assertion.

Si $B'$ est une seconde $A$-algèbre, $u:B\to B'$ un $A$-homomorphisme,
$\mathfrak{R}'$ un idéal de $B'$ tel que
$u(\mathfrak{R})\subset\mathfrak{R}'$ et $C'=B'/\mathfrak{R}'$ l'algèbre
quotient, on a un diagramme commutatif
\begin{equation}
\tag{20.5.11.3}
\label{0.20.5.11.3-fr}
\xymatrix{
  \mathfrak{R}/\mathfrak{R}^2 \ar[r]^-{\delta_{C/B/A}} \ar[d] &
  \Omega^1_{B/A}\otimes_B C \ar[d] \\
  \mathfrak{R}'/\mathfrak{R}'^2 \ar[r]_-{\delta_{C'/B'/A}} &
  \Omega^1_{B'/A}\otimes_{B'} C'
}
\end{equation}
où les flèches verticales proviennent de $u$ (20.5.2.4).
\end{env}

\begin{theorem}[20.5.12]
\label{0.20.5.12-fr}
Soient $A$ un anneau, $B$ une $A$-algèbre, $C$ la $A$-algèbre quotient
$B/\mathfrak{R}$, $u:B\to C$ l'homomorphisme canonique.
\begin{enumerate}
  \item[\rm(i)] On a une suite exacte de $C$-modules
\begin{equation}
\tag{20.5.12.1}
\label{0.20.5.12.1-fr}
  \mathfrak{R}/\mathfrak{R}^2
  \xrightarrow{\delta_{C/B/A}}\Omega^1_{B/A}\otimes_B C
  \xrightarrow{u_{C/B/A}}\Omega^1_{C/A}\longrightarrow 0,
\end{equation}
où $u_{C/B/A}$ et $\delta_{C/B/A}$ sont définis par (20.5.2.4) et
(20.5.11.2) respectivement.
  \item[\rm(ii)] Si l'on pose $E=B/\mathfrak{R}^2$, l'homomorphisme
canonique (20.5.2.4)
\[
  \Omega^1_{B/A}\otimes_B C\longrightarrow
  \Omega^1_{E/A}\otimes_E C
\]
est bijectif.
  \item[\rm(iii)] Les trois conditions suivantes sont équivalentes :
  \begin{enumerate}
    \item[\rm a)] $\delta_{C/B/A}$ est inversible à gauche.
    \item[\rm b)] Toute $A$-extension de $C$ par un $C$-module $L$, dont
l'image réciproque par $u:B\to C$ est $A$-triviale, est elle-même
$A$-triviale.
    \item[\rm c)] La $A$-algèbre $E=B/\mathfrak{R}^2$ est une extension
$A$-triviale de $C$ par $\mathfrak{R}/\mathfrak{R}^2$.
  \end{enumerate}
  \item[\rm(iv)] Il y a correspondance biunivoque canonique entre les
inverses à gauche de $\delta_{C/B/A}$ et les inverses à droite de
l'homomorphisme canonique $E\to C$.
\end{enumerate}
\end{theorem}

\begin{proof}
